%% ============================================================
%%  Chapter 1 — The Core Loop
%%  Production-Grade Agentic AI: LangGraph + Python Frameworks
%% ============================================================

\chapterquote{A good decision-making system is not one that never fails. It is
one that detects failure fast and corrects course before the damage
compounds.}{John Boyd, adapted}

A single LLM call can answer a question. It cannot own a goal.

That sentence describes the boundary between a language model application and an
agent system, and it is the boundary this entire book is concerned with crossing
responsibly. The moment a task requires a second step that depends on the result
of the first---retrieve a document, compare it to prior data, decide what to do
based on what you found---a one-shot model call is structurally insufficient. Not
because the model lacks reasoning quality, but because the surrounding system
lacks the orchestration to carry state, invoke tools, evaluate outcomes, and
decide what comes next.

This chapter builds the smallest complete agent loop you can put into production.
By the end, you will have a \code{StateGraph} with four typed nodes---Perceive,
Reason, Act, Reflect---wired by conditional edges, compiled with a checkpointer,
and instrumented with Langfuse spans on every node. Every design decision in that
loop is explained: why \code{TypedDict} coexists with Pydantic~v2, why reflection
is a typed control function rather than a narrative addition, why \code{.compile()}
is an architectural boundary and not a formality, and why cycle traces are
operational artifacts rather than optional debug output.

The running example throughout Part~I is an earnings-report agent: its task is to
retrieve a quarterly earnings report, compare flagged risk metrics to
prior-quarter thresholds, and produce a structured summary with escalation
signals. It is a genuine multi-step task---later steps depend on earlier results,
tool outputs shape plan revisions, and termination is not self-evident.

%% ------------------------------------------------------------
\section{AgentState: TypedDict + Pydantic~v2}
\label{sec:agentstate}
%% ------------------------------------------------------------

The state schema is the first thing you design, not the last. Before tools,
before prompts, before model selection, you define the data contract that all
nodes read and all nodes update. Everything else follows from it.

LangGraph requires state to be a \code{TypedDict}. This is not a
preference---it is a runtime constraint. LangGraph's graph executor, reducers,
checkpointer serialization, and streaming machinery all work over plain Python
dicts. Pydantic \code{BaseModel} instances carry methods, validators, and
metaclass machinery that conflict with this merge model.

The solution is to use both, at their correct boundaries. \code{AgentState} is a
\code{TypedDict} used internally by the graph at runtime. \code{AgentStateInput}
is a Pydantic~v2 \code{BaseModel} used as the validation gateway at the entry
point---when external code constructs an initial state before the first graph
invocation.

\begin{lstlisting}[language=Python,
  caption={\code{AgentState}: the runtime contract every node reads and writes},
  label={lst:agentstate}]
class AgentState(TypedDict):
    goal:             str                            # immutable; set at entry
    beliefs:          Annotated[list[str], add]      # append-only reducer
    intentions:       Annotated[list[str], add]
    token_budget:     int                            # hard ceiling, never None
    cycle_count:      int
    cycle_traces:     Annotated[list[dict], add]     # one entry per cycle
    messages:         Annotated[list[dict], add]
    last_action:      dict | None
    reflect_decision: str | None
    final_answer:     str | None
\end{lstlisting}

The \code{Annotated[list[str], add]} fields use LangGraph's built-in reducer.
When a node returns \code{{"beliefs": ["new observation"]}}, LangGraph appends
to the existing list rather than replacing it. \code{cycle\_traces} uses the
same reducer: every cycle appends one \code{CycleTrace} dict (\S\ref{sec:cycle-trace}).

The BDI mapping is intentional. \term{goal} is the desire---the objective the
agent pursues. \term{beliefs} is the belief store---the agent's accumulating
record of what it has observed. \term{intentions} is the intention
list---committed next steps derived from the current plan.

Pydantic validates only at the entry boundary, where external code constructs
the initial state. The validator below catches the two most common production
bugs before the first node ever runs.

\begin{lstlisting}[language=Python,
  caption={\code{AgentStateInput}: Pydantic guard at the graph entry point},
  label={lst:agentstateinput}]
class AgentStateInput(BaseModel):
    goal:         str
    token_budget: int
    model_config = {"frozen": True}

    @field_validator("token_budget")
    @classmethod
    def budget_must_be_positive(cls, v: int) -> int:
        if v <= 0:
            raise ValueError(f"token_budget must be > 0, got {v}")
        return v

    @field_validator("goal")
    @classmethod
    def goal_must_not_be_empty(cls, v: str) -> str:
        if not v.strip():
            raise ValueError("goal must not be empty")
        return v.strip()
\end{lstlisting}

\begin{warningbox}
\code{token\_budget} set to \code{None} does not raise a runtime error---it
silently propagates until a comparison against an integer fails deep inside a
node. The Pydantic validator here is cheap insurance that prevents a class of
production bugs that are extremely frustrating to trace. Never ship a graph entry
point without it.
\end{warningbox}

%% ------------------------------------------------------------
\section{PerceiveNode: Sensory Filtering and PII Redaction}
\label{sec:perceive}
%% ------------------------------------------------------------

Perception is the most under-engineered layer in most agent systems. PII that
enters \code{beliefs} or \code{messages} can leak into tool call arguments, audit
logs, and Langfuse traces. Malformed tool responses injected raw into the context
produce reasoning failures that look like model errors.

The Perceive node has three responsibilities: redact PII at ingestion, truncate
oversized inputs before they enter the context, and append a belief summarising
what was observed this cycle.

PII redaction uses Microsoft Presidio rather than regex patterns. Presidio's
\code{AnalyzerEngine} runs a pretrained NER model that generalises across names,
emails, phone numbers, national IDs, and financial identifiers without a pattern
library to maintain.

\begin{lstlisting}[language=Python,
  caption={PII redaction helpers used by \code{perceive\_node}},
  label={lst:perceive-helpers}]
_analyzer   = AnalyzerEngine()
_anonymizer = AnonymizerEngine()

def redact_pii(text: str, language: str = "en") -> tuple[str, int]:
    results = _analyzer.analyze(text=text, language=language)
    if not results:
        return text, 0
    anonymized = _anonymizer.anonymize(text=text,
                                       analyzer_results=results)
    return anonymized.text, len(results)

def _truncate_if_needed(text: str, max_chars: int = 4000) -> str:
    if len(text) <= max_chars:
        return text
    return (text[:max_chars]
            + f"\n[truncated -- {len(text) - max_chars} chars removed]")
\end{lstlisting}

The node itself is a pure \code{(state) -> dict} callable decorated with
\code{@observe} so Langfuse captures it as a named span. The decorator goes on
the node function, not on the graph runner---wrapping the runner would produce a
single opaque span for the entire run.

\begin{lstlisting}[language=Python,
  caption={\code{perceive\_node}: redact, truncate, increment counter, update beliefs},
  label={lst:perceive}]
@observe(name="perceive_node")
def perceive_node(state: AgentState) -> dict:
    updates: dict = {"cycle_count": state["cycle_count"] + 1}  # C1: must increment here
    redaction_count = 0
    messages = state.get("messages", [])

    if messages:
        last_msg = messages[-1]
        if last_msg.get("role") == "tool":
            raw       = last_msg.get("content", "")
            truncated = _truncate_if_needed(raw)
            redacted, n = redact_pii(truncated)
            redaction_count += n
            if redacted != raw:
                updates["_last_tool_content_redacted"] = redacted
        preview = str(last_msg.get("content", ""))[:120]
        updates["beliefs"] = [
            f"[cycle {state['cycle_count']}] "
            f"{last_msg.get('role','unknown')}: {preview}"
        ]

    langfuse_context.update_current_observation(metadata={
        "cycle_count":     state["cycle_count"] + 1,
        "redaction_count": redaction_count,
        "message_count":   len(messages),
    })
    return updates
\end{lstlisting}

\begin{warningbox}
\code{cycle\_count} must be incremented in \code{perceive\_node}, not in
\code{reflect\_node}. The counter tracks cycles \emph{entered}, not cycles
\emph{completed}. An incomplete cycle that consumed tokens still counts against
the \code{max\_cycles} limit. Omitting this increment means
\code{pre\_cycle\_gate} never sees the counter advance---making the
application-level guard permanently blind, even though \code{recursion\_limit}
will eventually fire. Both guards must work independently.
\end{warningbox}

\begin{tipbox}
Run \code{presidio-analyzer} in a separate process for high-throughput production
workloads. The NER model initialisation is expensive. The module-level
\code{\_analyzer} instance is sufficient for development and moderate loads; for
production consider serving Presidio as a sidecar.
\end{tipbox}

%% ------------------------------------------------------------
\section{ReasonNode: LiteLLM \texorpdfstring{$\to$}{->} Instructor
\texorpdfstring{$\to$}{->} Typed Union Output}
\label{sec:reason}
%% ------------------------------------------------------------

The Reason node is where the model call happens. But the model call is the middle
of the node, not the whole node. Before the call: assemble context, check the
token budget, route to the correct model. After the call: deduct tokens consumed,
parse the output into a typed object the orchestrator can act on mechanically.

LiteLLM provides the provider abstraction. Instructor patches the LiteLLM client
to add structured output enforcement via function calling or JSON mode, validates
the response against the target Pydantic schema, and retries automatically on
parse failure. The output is a discriminated union of three mutually exclusive
intent types.

\begin{lstlisting}[language=Python,
  caption={Discriminated union: the three intent types \code{reason\_node} can return},
  label={lst:reasonoutput}]
class ToolRequest(BaseModel):
    type:      Literal["tool_request"] = "tool_request"
    tool_name: str
    arguments: dict

class FinalAnswer(BaseModel):
    type:       Literal["final_answer"] = "final_answer"
    content:    str
    confidence: float   # 0.0--1.0

class ReasoningTrace(BaseModel):
    type:  Literal["reasoning_trace"] = "reasoning_trace"
    trace: str          # fed back into next Perceive cycle

ReasonOutput = ToolRequest | FinalAnswer | ReasoningTrace
\end{lstlisting}

\code{ReasoningTrace} is not an action. When the model needs to externalise
intermediate deliberation before committing to a tool call, it returns a
\code{ReasoningTrace}. The orchestrator feeds its \code{trace} content back into
the next Perceive cycle as a belief, then loops back through Reason again.
Nothing is executed.

\begin{lstlisting}[language=Python,
  caption={\code{reason\_node}: budget guard, model call, token deduction},
  label={lst:reason}]
class BudgetExhaustedError(RuntimeError):
    pass

_MINIMUM_SAFE_BUDGET = 500

@observe(name="reason_node")
def reason_node(state: AgentState) -> dict:
    remaining = state["token_budget"]
    if remaining < _MINIMUM_SAFE_BUDGET:
        raise BudgetExhaustedError(
            f"token_budget={remaining} below minimum safe "
            f"threshold of {_MINIMUM_SAFE_BUDGET}.")

    messages = [
        {"role": "system", "content": "You are a goal-directed "
         "reasoning engine. Output a ToolRequest, FinalAnswer, "
         "or ReasoningTrace."},
        *state.get("messages", [])[-10:],
        {"role": "user", "content":
            f"ACTIVE GOAL: {state['goal']}\n"
            f"CYCLE: {state['cycle_count']}\n"
            f"RECENT BELIEFS:\n"
            + "\n".join(state.get("beliefs", [])[-5:])
            + "\n\nREASON about the next best action."},
    ]
    output: ReasonOutput = client.chat.completions.create(
        model=state.get("_config_model", "gpt-4o-mini"),
        messages=messages,
        response_model=ReasonOutput,
        max_retries=2,
    )
    estimated_used = (len(str(messages)) + len(str(output))) // 4
    new_budget = max(0, remaining - estimated_used)
    langfuse_context.update_current_observation(
        usage={"input": estimated_used,
               "output": len(str(output)) // 4, "unit": "TOKENS"},
        metadata={"output_type": output.type,
                  "token_budget_remaining": new_budget})
    return {
        "messages":    [{"role": "assistant",
                         "content": output.model_dump_json()}],
        "last_action": output.model_dump(),
        "token_budget": new_budget,
    }
\end{lstlisting}

\begin{warningbox}
The token deduction above uses a character-count estimate. This is intentionally
approximate at this stage. Chapter~5 replaces it with LiteLLM's actual usage
response object and a proper callback-based accounting system. Never ship the
approximation to production.
\end{warningbox}

%% ------------------------------------------------------------
\section{ActionNode: Tool Dispatch and Audit}
\label{sec:action}
%% ------------------------------------------------------------

Every tool call is a commitment to a side effect that may be irreversible. The
\code{ToolInvocationRecord} emitted by this node is not optional
instrumentation---it is a structural output of every cycle that calls a tool.

\begin{notebox}
This chapter uses a hand-rolled \code{\_REGISTERED\_TOOLS} dict for clarity.
Chapter~14 replaces it with LangGraph's native \code{ToolNode}, which adds
parallel execution, \code{ToolException} routing, and \code{StructuredTool}
schema enforcement. The interface contract---\code{action\_node} receives a
\code{tool\_request} in \code{last\_action} and returns a \code{tool} message---is
identical in both implementations.
\end{notebox}

The audit record captures every invocation, including failures. Latency is
measured around the tool call only---not around the whole node---because tool
latency is the signal used to detect slow external dependencies.

\begin{lstlisting}[language=Python,
  caption={\code{ToolInvocationRecord}: mandatory fields on every tool call},
  label={lst:toolrecord}]
@dataclass
class ToolInvocationRecord:
    call:          str       # "<tool_name>(<json_args>)"
    result:        Any
    latency_ms:    float
    agent_id:      str
    tool_name:     str
    success:       bool
    error:         str | None = None
    timestamp_utc: str = field(
        default_factory=lambda:
            datetime.now(timezone.utc).isoformat()
    )
\end{lstlisting}

\begin{lstlisting}[language=Python,
  caption={\code{action\_node}: dispatch, time, record, return},
  label={lst:action}]
@observe(name="action_node")
def action_node(state: AgentState, config: dict) -> dict:
    last = state.get("last_action") or {}
    if last.get("type") != "tool_request":
        return {}
    tool_name = last["tool_name"]
    agent_id  = (config or {}).get("configurable", {}).get(
                  "agent_id", "unknown")
    if tool_name not in _REGISTERED_TOOLS:
        raise UnregisteredToolError(
            f"Model requested unregistered tool '{tool_name}'. "
            f"Available: {list(_REGISTERED_TOOLS)}")

    t0 = time.perf_counter_ns()
    try:
        result  = _REGISTERED_TOOLS[tool_name](**last["arguments"])
        success = True;  error = None
    except Exception as exc:
        result  = error = f"{type(exc).__name__}: {exc}"
        success = False
    latency_ms = (time.perf_counter_ns() - t0) / 1e6

    record = ToolInvocationRecord(
        call=f"{tool_name}({last['arguments']})", result=result,
        latency_ms=latency_ms, agent_id=agent_id,
        tool_name=tool_name, success=success, error=error)
    langfuse_context.update_current_observation(metadata={
        "tool_name": tool_name, "success": success,
        "latency_ms": latency_ms, "error": error})
    return {
        "messages":    [{"role": "tool", "content": str(result),
                         "tool_call_id": tool_name}],
        "last_action": {**last, **record.__dict__},
    }
\end{lstlisting}

%% ------------------------------------------------------------
\section{ReflectNode: Enum-Gated Typed Decision}
\label{sec:reflect}
%% ------------------------------------------------------------

Reflection is a control function, not a narrative addition to the prompt. The
five enum values produced by this node are the complete, closed decision space
for the loop transition after every action.

\begin{itemize}
  \item \code{CONTINUE} --- advance to the next cycle normally
  \item \code{REVISE} --- update the goal or plan before the next cycle
  \item \code{RETRY} --- re-execute the same tool with no plan change
  \item \code{TERMINATE} --- task is complete; set \code{final\_answer} and exit
  \item \code{ESCALATE} --- task cannot be completed; surface to human
\end{itemize}

The distinction between \code{REVISE} and \code{RETRY} is operationally
important. \code{RETRY} means the plan is correct but the execution
failed---try again. \code{REVISE} means the execution succeeded but revealed
that the current plan is wrong---change the plan.

\begin{lstlisting}[language=Python,
  caption={\code{ReflectionOutput}: the typed contract driving every routing decision},
  label={lst:reflectionoutput}]
class ReflectionOutput(BaseModel):
    decision:          Literal[
                         "CONTINUE","REVISE","RETRY",
                         "TERMINATE","ESCALATE"]
    goal_delta:        Literal["ADVANCED","NO_CHANGE","REGRESSED"]
    technical_success: bool
    reason_code:       str
    summary:           str
\end{lstlisting}

\begin{lstlisting}[language=Python,
  caption={\code{reflect\_node}: evaluate, decide, optionally set \code{final\_answer}},
  label={lst:reflect}]
@observe(name="reflect_node")
def reflect_node(state: AgentState, config: dict) -> dict:
    last  = state.get("last_action") or {}
    model = ((config or {}).get("configurable", {})
              .get("reflect_model", "gpt-4o-mini"))
    output: ReflectionOutput = client.chat.completions.create(
        model=model,
        messages=[
            {"role": "system",
             "content": "Choose: CONTINUE|REVISE|RETRY|"
                        "TERMINATE|ESCALATE."},
            {"role": "user", "content":
                f"GOAL: {state['goal']}\n"
                f"CYCLE: {state['cycle_count']}\n"
                f"ACTION: {last.get('call','none')}\n"
                f"SUCCESS: {last.get('success',False)}\n"
                f"ERROR: {last.get('error') or 'none'}\n"
                f"RESULT: {str(last.get('result',''))[:500]}"},
        ],
        response_model=ReflectionOutput,
        max_retries=2,
    )
    langfuse_context.update_current_observation(metadata={
        "decision": output.decision, "goal_delta": output.goal_delta,
        "reason_code": output.reason_code})
    updates: dict = {
        "reflect_decision": output.decision,
        "beliefs": [f"[reflect cycle={state['cycle_count']}] "
                    f"{output.summary}"],
    }
    if output.decision == "TERMINATE":
        updates["final_answer"] = (
            state.get("final_answer") or next(
                (m.get("content","") for m in
                 reversed(state.get("messages",[]))
                 if m.get("role") == "assistant"),
                "Task complete."))
    return updates
\end{lstlisting}

\begin{notebox}
The five \code{decision} values are locked. Chapter~2's routing function maps
each value to a graph branch. Any additional value added to this enum without
updating the routing branches will hit the \code{else} clause and silently
terminate the run.
\end{notebox}

%% ------------------------------------------------------------
\section{StateGraph Wiring and \texttt{.compile()} as an
         Architectural Boundary}
\label{sec:graph}
%% ------------------------------------------------------------

\code{.compile()} is a named architectural boundary for three reasons that matter
operationally. First, it validates graph structure at compile time: misconfigured
edges are caught before the graph runs a single cycle. Second, it enables
streaming via \code{graph.stream()}. Third, it connects the checkpointer: every
node execution is automatically persisted to the configured store.

The two routing functions encode the complete control flow. \code{route\_from\_reason}
acts on the \code{type} field of \code{last\_action}. \code{route\_from\_reflect}
acts on the five locked \code{reflect\_decision} values. \code{TERMINATE} and
\code{ESCALATE} both route to \code{END} here; Chapter~2 introduces a dedicated
\code{escalation\_node} that separates the two paths so that human-handoff
telemetry is distinct from normal completion telemetry.

\begin{lstlisting}[language=Python,
  caption={Routing functions and \code{StateGraph} wiring},
  label={lst:graph}]
def route_from_reason(state: AgentState) -> str:
    t = (state.get("last_action") or {}).get("type", "")
    if t == "final_answer":    return "end"
    if t == "reasoning_trace": return "perceive_node"
    return "action_node"

def route_from_reflect(state: AgentState) -> str:
    d = state.get("reflect_decision") or "CONTINUE"
    return {"CONTINUE":  "perceive_node",
            "REVISE":    "perceive_node",
            "RETRY":     "action_node",
            "TERMINATE": "end",
            "ESCALATE":  "end"}.get(d, "end")

def build_graph(policy: TerminationPolicy,
                checkpointer=None):
    b = StateGraph(AgentState)
    for name, fn in [("perceive_node", perceive_node),
                     ("reason_node",   reason_node),
                     ("action_node",   action_node),
                     ("reflect_node",  reflect_node)]:
        b.add_node(name, fn)
    b.set_entry_point("perceive_node")
    b.add_edge("perceive_node", "reason_node")
    b.add_conditional_edges("reason_node", route_from_reason,
        {"action_node": "action_node",
         "perceive_node": "perceive_node", "end": END})
    b.add_edge("action_node", "reflect_node")
    b.add_conditional_edges("reflect_node", route_from_reflect,
        {"perceive_node": "perceive_node",
         "action_node": "action_node", "end": END})
    return b.compile(checkpointer=checkpointer,
                     recursion_limit=policy.max_cycles * 4)
\end{lstlisting}

\begin{notebox}
\code{recursion\_limit} is derived from \code{policy.max\_cycles} at compile
time---not from a hard-coded constant. This means the structural ceiling and the
application ceiling are always proportional: a policy change updates both without
requiring a separate edit to \code{compile()}.
\end{notebox}

\code{recursion\_limit=policy.max\_cycles * 4} is a hard guard: LangGraph raises
\code{GraphRecursionError} when it is reached. Chapter~2 explains why increasing
this number is not a valid alternative to proper termination logic.

%% ------------------------------------------------------------
\section{Embedded vs.\ Explicit Reflection}
\label{sec:reflection-modes}
%% ------------------------------------------------------------

Embedded reflection means the next \code{ReasonNode} call receives the last
\code{ToolMessage} in context and implicitly decides what to do next through its
regular reasoning pass---no separate model call, no separate node. Embedded
reflection costs less: one fewer model call per cycle, lower latency, one fewer
span in the trace tree.

The promotion criteria for moving to explicit \code{ReflectNode} are four
specific production requirements:

\begin{enumerate}
  \item Independent audit of self-evaluation is required by policy.
  \item Action safety requires a separate gate before any information that could
        influence the next action is mixed with new context.
  \item Recovery logic is becoming opaque inside the Reason prompt.
  \item A different, smaller model for reflection is economically justified.
\end{enumerate}

Start with \code{"explicit"} for any task that involves external tool calls. Use
\code{"embedded"} only for tasks with no side effects, short horizons, and no
audit requirements.

%% ------------------------------------------------------------
\section{Langfuse Span Emission via \texttt{@observe}}
\label{sec:observe}
%% ------------------------------------------------------------

\code{@observe} must decorate the \textbf{node function}, not the graph runner or
\code{graph.invoke()}. Wrapping the graph runner produces one Langfuse span for
the entire run. Decorating each node function produces a trace tree: four child
spans per cycle---\code{perceive\_node}, \code{reason\_node},
\code{action\_node}, \code{reflect\_node}---each with its own duration, token
usage, and metadata.

The \code{name=} parameter on \code{@observe} must exactly match the string used
in \code{builder.add\_node()}. Langfuse Studio correlates span names to graph
node names when they are identical, enabling trace-to-graph overlays.

Session grouping is handled by passing \code{session\_id} through
\code{config["configurable"]}:

\begin{lstlisting}[language=Python,
  caption={Passing session and model config through \code{graph.invoke()}},
  label={lst:session}]
result = graph.invoke(
    initial_state,
    config={
        "configurable": {
            "agent_id":      "agent-001",
            "session_id":    "session-abc123",
            "model":         "gpt-4o",
            "reflect_model": "gpt-4o-mini",
        }
    },
)
\end{lstlisting}

%% ------------------------------------------------------------
\section{Cycle Trace Record: 13 Mandatory Fields}
\label{sec:cycle-trace}
%% ------------------------------------------------------------

A Langfuse span is push telemetry to an external service. A \code{CycleTrace} is
a typed Python dataclass appended to \code{AgentState.cycle\_traces} after every
cycle. Chapter~2's infinite loop and goal divergence detectors read this
field every cycle---it must be present or those detectors have nothing to act on.

\begin{lstlisting}[language=Python,
  caption={\code{CycleTrace}: 13 mandatory fields capturing one complete cycle},
  label={lst:cycletrace}]
@dataclass
class TokenUsage:
    input:            int
    output:           int
    cumulative_total: int

@dataclass
class CycleTrace:
    cycle_number:             int
    goal_summary:             str
    context_package_summary:  str
    reason_output_type:  Literal["TOOL_CALL","FINAL_ANSWER",
                                  "REASONING_TRACE"]
    action_invoked:      str | None
    validation_result:   Literal["PASSED","REJECTED","SKIPPED"]
    tool_result_summary: str | None
    reflection_decision: Literal["CONTINUE","REVISE","RETRY",
                                  "TERMINATE","ESCALATE"]
    goal_delta:     Literal["ADVANCED","NO_CHANGE","REGRESSED"]
    termination_status: Literal["RUNNING","COMPLETED",
                                 "RESOURCE_LIMIT","FAILURE",
                                 "ESCALATED"]
    token_usage:    TokenUsage
    wall_time_ms:   float
    reason_code:    str
\end{lstlisting}

\code{CycleTrace} instances are written by a thin helper called after
\code{reflect\_node} completes. Because \code{cycle\_traces} uses the
\code{Annotated[list[dict], add]} reducer, writing it is a single return value:

\begin{lstlisting}[language=Python,
  caption={Writing \code{CycleTrace} back into state after each cycle},
  label={lst:cycletrace-write}]
def record_cycle_trace(state: AgentState,
                       trace: CycleTrace) -> dict:
    """Call from a post-reflect wrapper or orchestrator."""
    return {"cycle_traces": [asdict(trace)]}
\end{lstlisting}

\code{termination\_status} uses different vocabulary from
\code{reflect\_decision}. A \code{TERMINATE} decision maps to \code{COMPLETED}.
A task stopped by the \code{recursion\_limit} guard maps to
\code{RESOURCE\_LIMIT}. Aggregating them under a single ``stopped'' label hides a
significant quality signal.

%% ------------------------------------------------------------
\section{pytest Suite: Node Isolation Tests}
\label{sec:tests}
%% ------------------------------------------------------------

The node functions are pure \code{(state: AgentState, config: dict) -> dict}
callables. No graph compilation is required to test them. The shared fixture
constructs a minimal valid \code{AgentState} for every test.

\begin{lstlisting}[language=Python,
  caption={\code{make\_test\_state}: shared fixture for all node tests},
  label={lst:conftest}]
# tests/conftest.py
def make_test_state(**overrides) -> AgentState:
    base: AgentState = {
        "goal": "Test goal",     "beliefs": [],
        "intentions": [],        "token_budget": 4000,
        "cycle_count": 1,        "cycle_traces": [],
        "messages": [],          "last_action": None,
        "reflect_decision": None,"final_answer": None,
    }
    base.update(overrides)
    return base
\end{lstlisting}

Each node is tested in isolation. The key assertions express the
\emph{contract}---what the node must guarantee---not its implementation details.

\begin{lstlisting}[language=Python,
  caption={Node isolation tests: one assertion per contract guarantee},
  label={lst:tests}]
# LANGFUSE_ENABLED=false disables SDK network calls for the test process
os.environ.setdefault("LANGFUSE_ENABLED", "false")

# --- perceive_node: cycle_count must be incremented ---
def test_cycle_count_incremented():
    state  = make_test_state(cycle_count=3)
    result = perceive_node(state)
    assert result["cycle_count"] == 4

# --- perceive_node: PII must be redacted before beliefs are written ---
def test_pii_is_redacted():
    state = make_test_state(messages=[
        {"role": "tool", "content": "Email: john@example.com"}])
    with patch("src.agent.nodes.perceive._analyzer") as ma, \
         patch("src.agent.nodes.perceive._anonymizer") as mb:
        mb.anonymize.return_value.text = "Email: <EMAIL_ADDRESS>"
        ma.analyze.return_value = [MagicMock()]
        result = perceive_node(state)
    assert "<EMAIL_ADDRESS>" in result["_last_tool_content_redacted"]

# --- reason_node: exhausted budget must raise, never silently continue ---
def test_budget_exhausted_raises():
    with pytest.raises(BudgetExhaustedError):
        reason_node(make_test_state(token_budget=100))

# --- action_node: every invocation must produce the four audit fields ---
def test_record_mandatory_fields():
    register_tool("echo", lambda msg: f"echo: {msg}")
    state = make_test_state(last_action={
        "type": "tool_request", "tool_name": "echo",
        "arguments": {"msg": "hello"}})
    result = action_node(state, {"configurable": {"agent_id": "t"}})
    assert all(k in result["last_action"]
               for k in ["call", "result", "latency_ms", "agent_id"])

# --- reflect_node: TERMINATE decision must write final_answer ---
def test_terminate_sets_final_answer():
    state = make_test_state(
        messages=[{"role": "assistant", "content": "Done."}],
        last_action={"call": "f()", "result": "ok", "success": True})
    mock_out = ReflectionOutput(
        decision="TERMINATE", goal_delta="ADVANCED",
        technical_success=True, reason_code="goal_met",
        summary="Goal met.")
    with patch("src.agent.nodes.reflect.client") as mc:
        mc.chat.completions.create.return_value = mock_out
        result = reflect_node(state, {})
    assert result["reflect_decision"] == "TERMINATE"
    assert result.get("final_answer") is not None
\end{lstlisting}

Each node is tested without a compiled graph and without Langfuse network calls.
Chapter~28 extends this suite to graph-level integration tests that compile the
full graph, use \code{SqliteSaver} checkpointing, and assert on multi-cycle state
evolution.

\begin{takeawaybox}
The core loop established in this chapter is a fixed contract for the rest of the
book. \code{AgentState} is a \code{TypedDict}; Pydantic validates at the
boundary only. \code{PerceiveNode} redacts PII via Presidio before any content
reaches the model. \code{ReasonNode} returns a discriminated union the
orchestrator routes on mechanically. \code{ActionNode} emits a
\code{ToolInvocationRecord} on every execution including failures.
\code{ReflectNode} produces five locked enum values. \code{.compile()} validates
structure, enables streaming, and connects the checkpointer. \code{@observe}
decorates node functions, not graph runners. \code{CycleTrace} is appended to
state after every cycle---Chapter~2's detectors depend on it. Chapter~2 asks the
question this chapter left open: what happens when the loop does not stop?
\end{takeawaybox}
